% !TeX root = ../main.tex
\chapter{Problem Background and Definition}
\label{chap:detailed-design}

\section{Purpose and Scope}
This chapter documents the hardware and software design of the
Multi-Rail Isolated Power Supply Monitoring System. The design monitors six
power-supply rails, measures rail voltage and load current, preserves electrical
isolation between floating supplies and the controller, evaluates measurements
against user-configurable alarm limits, presents local status on the controller
display, and publishes measurements and status to an EPICS control-system
interface.

This design is intended to establish the major analog and digital
functional blocks, document supporting engineering calculations, define the
hardware schematic baseline, and provide software flowcharts detailed enough to
support implementation.

\section{Problem Definition}
The project addresses the need to monitor several low-voltage power-supply rails
that do not necessarily share a common reference. A conventional
monitoring circuit that directly connects the grounds together
would defeat the electrical isolation of the supplies and could create unwanted
return paths. However, the user needs a single location from which to observe rail voltage, rail current, alarm state, communication status, and the overall condition of the power system.

The required monitoring system must therefore make accurate measurements on
six separate rails while preserving isolation between the monitored domains. It
must support both positive and negative rails, detect sensor or communication
faults, provide configurable voltage and current limits, retain configuration
information after reset, present useful local status, and make the same data
available to the facility control system over Ethernet. The monitored rails in
the present design are +5 V A, +5 V B, +/-5 V, and +/-15 V.

\section{Proposed Solution}
The proposed solution assigns one LTC2945 voltage/current monitor to each power
rail. Load current is converted to a small differential voltage by a precision
high-side shunt resistor and the LTC2945 digitizes both the shunt voltage and the
rail voltage. Each sensor remains referenced to its own monitored supply domain.
The I\textsuperscript{2}C communication path is then isolated so that digital
measurement data can be transferred to a common Arduino GIGA R1 WiFi controller
without directly bonding the isolated supply returns together.

The Arduino GIGA performs channel acquisition, engineering-unit conversion,
negative-rail sign handling, alarm evaluation, configuration storage, display
updates, and network communication. A local GIGA Display Shield provides
immediate status at the hardware. Ethernet transports measurement and status
data using the PSC protocol to an EPICS IOC, where PSCDriver and EPICS records
make the values available to a CS-Studio Phoebus operator interface. This
architecture separates the rail-side analog measurement functions from the
controller and supervisory-control functions while maintaining a single user
interface for the complete system.